% GENERATED FILE. Edit canonical lessons, then run npm run generate:books.
% canonical-source-hash: fnv1a64:e3b1b1fc290c25ee
% canonical-lessons: FA-C345-qesse, FA-C345-ruzi-ruzegari, FA-C345-yeki-bud-yeki-nabud, FA-C345-qesse-gu, FA-C345-ravi

\chapter{Once upon a time}
\label{ch:fa-once-upon-a-time}
\hlchaptermodality{345}

\begin{chapteropening}
\textbf{By the end of this chapter:} \emph{I can begin and tell a folk tale, and talk about the storyteller and the narrator.}

Complete the last lesson of chapter 345: I can begin and tell a folk tale, and talk about the storyteller and the narrator.
\end{chapteropening}

\section[\texorpdfstring{\fa{قصه} (qesse)}{قصه (qesse)}]{\fa{قصه} (qesse) — a tale, a story}

\label{lesson:FA-C345-qesse}



\begin{quote}
Before the new one: say the Persian for saffron rice pudding, then the Persian for rosewater.
\end{quote}

\subsection*{You'll want to know: \fa{قصه}}
\textbf{\fa{قصه}} (\emph{qesse}) — \textquotedblleft{}a tale, a story\textquotedblright{}.

\begin{quote}
\fa{مادربزرگم} \fa{هر} \fa{شب} \fa{برای} \fa{ما} \fa{قصه} \fa{میگفت}.
\end{quote}

\begin{quote}
\emph{Mâdarbozorgam har shab barâ-ye mâ qesse migoft.}
\end{quote}

\begin{quote}
Every night my grandmother used to tell us a story.
\end{quote}

Close to \emph{dâstân}, but often a folk tale told aloud. \emph{Qesse goftan} \textquotedblleft{}to tell a tale\textquotedblright{}.

\begin{cousinweb}
\emph{qesse} is borrowed from Arabic \emph{qiṣṣa}, \textquotedblleft{}story\textquotedblright{} (EtymDB).
\end{cousinweb}

\begin{grammarlens}[title={What you've built}]
One more word for telling, explaining and describing at B1.
\end{grammarlens}

\subsection*{Your turn}
\begin{itemize}
\raggedright
  \item \emph{Say it:} \emph{qesse}
  \item \emph{Say it:} \emph{qesse}, once more
  \item \emph{Say it:} a sentence of your own with \emph{qesse}
  \item \emph{Recall:} say the Persian for saffron rice pudding, then the Persian for rosewater, then say \emph{qesse} again
\end{itemize}

\begin{tcolorbox}[breakable,colback=teal!4,colframe=teal!35!black,title={Before you move on}]
What does \fa{قصه} mean? (\textquotedblleft{}A tale, a story\textquotedblright{}.) Say it once more.
\end{tcolorbox}
\section[\texorpdfstring{\fa{روزی} \fa{روزگاری} (ruzi ruzegâri)}{روزی روزگاری (ruzi ruzegâri)}]{\fa{روزی} \fa{روزگاری} (ruzi ruzegâri) — once upon a time}

\label{lesson:FA-C345-ruzi-ruzegari}



\begin{quote}
Before the new one: say the Persian for rosewater, then the Persian for a tale.
\end{quote}

\subsection*{You'll want to know: \fa{روزی} \fa{روزگاری}}
\textbf{\fa{روزی} \fa{روزگاری}} (\emph{ruzi ruzegâri}) — \textquotedblleft{}once upon a time\textquotedblright{}.

\begin{quote}
\fa{روزی} \fa{روزگاری} \fa{در} \fa{یک} \fa{روستای} \fa{کوچک} \fa{پیرزنی} \fa{زندگی} \fa{میکرد}.
\end{quote}

\begin{quote}
\emph{Ruzi ruzegâri dar yek rustâ-ye kuchak pirzani zendegi mikard.}
\end{quote}

\begin{quote}
Once upon a time an old woman lived in a little village.
\end{quote}

Literally \textquotedblleft{}one day, in one time\textquotedblright{}. \emph{ruzegâr} \textquotedblleft{}time, age\textquotedblright{}. The past continuous sets the scene.

\begin{cousinweb}
\emph{ruzi ruzegâri} repeats \emph{ruz} \textquotedblleft{}day\textquotedblright{} in \emph{ruzi} \textquotedblleft{}one day\textquotedblright{} and \emph{ruzegâri} \textquotedblleft{}in one age\textquotedblright{}, from \emph{ruzegâr} \textquotedblleft{}time, age\textquotedblright{}. EtymDB takes \emph{ruz} back through Middle Persian to Old Persian \emph{rauča-}.
\end{cousinweb}

\begin{grammarlens}[title={What you've built}]
One more word for telling, explaining and describing at B1.
\end{grammarlens}

\subsection*{Your turn}
\begin{itemize}
\raggedright
  \item \emph{Say it:} \emph{ruzi ruzegâri}
  \item \emph{Say it:} \emph{ruzi ruzegâri}, once more
  \item \emph{Say it:} a sentence of your own with \emph{ruzi ruzegâri}
  \item \emph{Recall:} say the Persian for rosewater, then the Persian for a tale, then say \emph{ruzi ruzegâri} again
\end{itemize}

\begin{tcolorbox}[breakable,colback=teal!4,colframe=teal!35!black,title={Before you move on}]
What does \fa{روزی} \fa{روزگاری} mean? (\textquotedblleft{}Once upon a time\textquotedblright{}.) Say it once more.
\end{tcolorbox}
\section[\texorpdfstring{\fa{یکی} \fa{بود} \fa{یکی} \fa{نبود} (yeki bud yeki nabud)}{یکی بود یکی نبود (yeki bud yeki nabud)}]{\fa{یکی} \fa{بود} \fa{یکی} \fa{نبود} (yeki bud yeki nabud) — there was one and there was none (the opening of a folk tale)}

\label{lesson:FA-C345-yeki-bud-yeki-nabud}



\begin{quote}
Before the new one: say the Persian for a tale, then the Persian for once upon a time.
\end{quote}

\subsection*{You'll want to know: \fa{یکی} \fa{بود} \fa{یکی} \fa{نبود}}
\textbf{\fa{یکی} \fa{بود} \fa{یکی} \fa{نبود}} (\emph{yeki bud yeki nabud}) — \textquotedblleft{}there was one and there was none (the opening of a folk tale)\textquotedblright{}.

\begin{quote}
\fa{قصه} \fa{با} \fa{یکی} \fa{بود} \fa{یکی} \fa{نبود} \fa{شروع} \fa{شد}.
\end{quote}

\begin{quote}
\emph{Qesse bâ yeki bud yeki nabud shoru' shod.}
\end{quote}

\begin{quote}
The tale began with There was one, there was none.
\end{quote}

The traditional opening of a Persian folk tale, often followed by \emph{gheyr az khodâ hich-kas nabud} \textquotedblleft{}and apart from God there was nobody\textquotedblright{}.

\begin{cousinweb}
\emph{yeki bud yeki nabud} is built from native Persian words only: \emph{yeki} \textquotedblleft{}one\textquotedblright{}, \emph{bud} \textquotedblleft{}was\textquotedblright{} and \emph{nabud} \textquotedblleft{}was not\textquotedblright{}.
\end{cousinweb}

\begin{grammarlens}[title={What you've built}]
One more word for telling, explaining and describing at B1.
\end{grammarlens}

\subsection*{Your turn}
\begin{itemize}
\raggedright
  \item \emph{Say it:} \emph{yeki bud yeki nabud}
  \item \emph{Say it:} \emph{yeki bud yeki nabud}, once more
  \item \emph{Say it:} a sentence of your own with \emph{yeki bud yeki nabud}
  \item \emph{Recall:} say the Persian for a tale, then the Persian for once upon a time, then say \emph{yeki bud yeki nabud} again
\end{itemize}

\begin{tcolorbox}[breakable,colback=teal!4,colframe=teal!35!black,title={Before you move on}]
What does \fa{یکی} \fa{بود} \fa{یکی} \fa{نبود} mean? (\textquotedblleft{}There was one and there was none (the opening of a folk tale)\textquotedblright{}.) Say it once more.
\end{tcolorbox}
\section[\texorpdfstring{\fa{قصه} \fa{گو} (qesse gu)}{قصه گو (qesse gu)}]{\fa{قصه} \fa{گو} (qesse gu) — a storyteller}

\label{lesson:FA-C345-qesse-gu}



\begin{quote}
Before the new one: say the Persian for once upon a time, then the Persian for there was one and there was none.
\end{quote}

\subsection*{You'll want to know: \fa{قصه} \fa{گو}}
\textbf{\fa{قصه} \fa{گو}} (\emph{qesse gu}) — \textquotedblleft{}a storyteller\textquotedblright{}.

\begin{quote}
\fa{قصه} \fa{گو} \fa{داستان} \fa{را} \fa{با} \fa{صدای} \fa{بلند} \fa{تعریف} \fa{کرد}.
\end{quote}

\begin{quote}
\emph{Qesse gu dâstân râ bâ sedâ-ye boland ta'rif kard.}
\end{quote}

\begin{quote}
The storyteller told the story in a loud voice.
\end{quote}

\emph{qesse} \textquotedblleft{}tale\textquotedblright{} + \emph{gu}, the present stem of \emph{goftan} \textquotedblleft{}to say\textquotedblright{}.

\begin{cousinweb}
\emph{qesse} \textquotedblleft{}tale\textquotedblright{} is borrowed from Arabic \emph{qiṣṣa} (EtymDB), and \emph{gu-} is the present stem of the native verb \emph{goftan} \textquotedblleft{}to say\textquotedblright{}: a tale-sayer.
\end{cousinweb}

\begin{grammarlens}[title={What you've built}]
One more word for telling, explaining and describing at B1.
\end{grammarlens}

\subsection*{Your turn}
\begin{itemize}
\raggedright
  \item \emph{Say it:} \emph{qesse gu}
  \item \emph{Say it:} \emph{qesse gu}, once more
  \item \emph{Say it:} a sentence of your own with \emph{qesse gu}
  \item \emph{Recall:} say the Persian for once upon a time, then the Persian for there was one and there was none, then say \emph{qesse gu} again
\end{itemize}

\begin{tcolorbox}[breakable,colback=teal!4,colframe=teal!35!black,title={Before you move on}]
What does \fa{قصه} \fa{گو} mean? (\textquotedblleft{}A storyteller\textquotedblright{}.) Say it once more.
\end{tcolorbox}
\section[\texorpdfstring{\fa{راوی} (râvi)}{راوی (râvi)}]{\fa{راوی} (râvi) — a narrator}

\label{lesson:FA-C345-ravi}



\begin{quote}
Before the new one: say the Persian for there was one and there was none, then the Persian for a storyteller.
\end{quote}

\subsection*{You'll want to know: \fa{راوی}}
\textbf{\fa{راوی}} (\emph{râvi}) — \textquotedblleft{}a narrator\textquotedblright{}.

\begin{quote}
\fa{در} \fa{این} \fa{رمان} \fa{راوی} \fa{یک} \fa{بچه} \fa{است}.
\end{quote}

\begin{quote}
\emph{Dar in român râvi yek bachche ast.}
\end{quote}

\begin{quote}
In this novel the narrator is a child.
\end{quote}

Plural \emph{râviyân}. It also names someone who passes on a traditional report.

\begin{cousinweb}
\emph{râvi} is usually traced to Arabic \emph{rāwī}, \textquotedblleft{}narrator, one who passes on a report\textquotedblright{}; EtymDB lists no entry.
\end{cousinweb}

\begin{grammarlens}[title={What you've built}]
One more word for telling, explaining and describing at B1.
\end{grammarlens}

\subsection*{Your turn}
\begin{itemize}
\raggedright
  \item \emph{Say it:} \emph{râvi}
  \item \emph{Say it:} \emph{râvi}, once more
  \item \emph{Say it:} a sentence of your own with \emph{râvi}
  \item \emph{Recall:} say the Persian for there was one and there was none, then the Persian for a storyteller, then say \emph{râvi} again
\end{itemize}

\begin{tcolorbox}[breakable,colback=teal!4,colframe=teal!35!black,title={Before you move on}]
What does \fa{راوی} mean? (\textquotedblleft{}A narrator\textquotedblright{}.) Say it once more.
\end{tcolorbox}
